\section{The unsigned positional comparison}
\label{p14:positions:section}

The scalar approximation sums over the locations of the starts. A version
of the same stopped expansion keeps their unsigned positions on the
original integer grid. Its reference consists of single Poisson arrivals
at sites and double arrivals along good arithmetic pairs. The statement
below retains the $4/5$ error established by this vector expansion; the
optimal scalar $4/3$ estimate does not by itself improve that rate.

Put $G=G_N$ and $p=2^{-L}$. The good edges are those of
Section~\ref{p1:arith:section} with both endpoints in $G$. For
$\boldsymbol z\in\N^{I_N}$, write $|\boldsymbol z|=\sum_xz_x$.

Let $E_N$ be the good-edge set on $G$,
$q_e=\P(J_x=J_y=1)$ for $e=\{x,y\}$,
$\beta=\beta_N=\sum_e q_e$, and define
\begin{equation}
 \begin{gathered}
 \iota_x=\sum_{e\ni x}q_e,\qquad \nu_x=p-\iota_x,\\
 U_x\sim\Pois(\nu_x),\qquad V_e\sim\Pois(q_e),\qquad
 H_x=U_x+\sum_{e\ni x}V_e,
 \end{gathered}
 \label{p1:limits:field-definition}
\end{equation}
with all the $U_x,V_e$ independent.  Here and below $\iota_x=0$
outside $G$, and $U_x$ is also defined at those starts.
The good-edge degree bound and the bound for $q_e$ imply
$\iota_x/p\ll Bp2^{L/3}=o(1)$, so the intensities $\nu_x$ are
positive for large $N$.  Write $\mathsf H_N=\Law((H_x)_{x\in I_N})$
and $\mathsf U_N=\bigotimes_{x\in I_N}\Pois(\nu_x)$.
Their total-count laws are respectively
\begin{equation}
 \begin{gathered}
 H_{\lambda,\theta}
   :=\Law\bigl(\Pois(\lambda-\theta)
                     +2\Pois(\theta/2)\bigr),\qquad
 \Pois(\lambda-\theta),\\
 \theta=2\beta,\qquad\lambda=Np.
 \end{gathered}
 \label{p1:limits:compound-target}
\end{equation}
where the two Poisson variables in the first law are independent.

\paragraph{Exact identities of the reference field.}
Every $H_x$ has law $\Pois(p)$, and all its polynomial moments are
finite. Its variance is $p$; for distinct sites its covariance is
$q_{\{x,y\}}$ on a retained edge and zero otherwise. Two disjoint
finite coordinate blocks are independent exactly when no positive-rate
retained edge crosses between them. The exact bivariate law and the
proof of this full block statement are given in
Section~\ref{p14:comp:positional-proof}.
These identities concern $H$, not the arithmetic indicators.
For example, a retained regular arithmetic pair has covariance $q_e-p^2$
and a regular indicator has variance $p(1-p)$.
The nonnegative compensation $\nu_x=p-\iota_x$ is part of the
reference construction; it does not decompose the actual indicators
into independent arrivals.

\begin{lemma}[Unsigned positional compound-Poisson comparison]
\label{p1:limits:vector-reference}
For every fixed integer $j\ge0$ and $\varepsilon>0$, uniformly
in the critical band,
\begin{equation}
 \sum_{\boldsymbol z\in\mathbb N^{I_N}}
 (1+|\boldsymbol z|)^j
 \big|\P((J_x)_x=\boldsymbol z)-\mathsf H_N\{\boldsymbol z\}\big|
       \ll_{C,j,\varepsilon}N^{-4/5+\varepsilon},
 \qquad |\boldsymbol z|=\sum_xz_x.
 \label{p1:limits:vector-weighted}
\end{equation}
\end{lemma}

The series is summable. One threshold is chosen after $C,j,\varepsilon$
and before $N,L$, the probability space and all independent Poisson
arrivals with the specified rates. A finite product of Poisson measures
provides such a realization, and every realization gives the same law.
The proof is in Section~\ref{p14:comp:positional-proof}. It retains the
coordinates in the finite convolution expansion, bounds the same positive
boundary and independent deficits, and uses the independently established
fixed moments to justify the polynomial weights.

This field keeps the integer locations of both endpoints of each pair.
It does not record run signs or exact excess lengths. The individual
weights $q_e$ specify the target: the scalar second factorial cumulant
alone does not determine their spatial arrangement. The companion also
uses this comparison to justify pair deletion and to calculate the
remaining covariance and signed Palm terms. Those conclusions do not
assert an exponential-moment bound for the original count.
